% 02 - SYSTEME ET PERIMETRE
% ==============================================================
\sectionintro{02 · Réalisation}{Système et périmètre implémenté}{Un bloc fonctionnel unique porte la logique métier ; le programme principal reste volontairement léger.}

\twocolpage{
  \subsection{Structure du projet XAE}
  \begin{tikzpicture}[node distance=6.2mm,font=\scriptsize]
    \node[draw=ekprimary!45,fill=eksoft,rounded corners=2mm,minimum width=53mm,minimum height=9mm,align=left] (root) {\textbf{TQS\_Test\_CHATER}};
    \node[draw=ekline,fill=white,rounded corners=2mm,below=of root,minimum width=53mm,minimum height=8mm,align=left] (plc) {PLC · projet PLC};
    \node[draw=ekline,fill=white,rounded corners=2mm,below=of plc,minimum width=53mm,minimum height=8mm,align=left] (dut) {DUTs · \texttt{E\_ModeCarrousel}\\\phantom{DUTs · }\texttt{E\_EtatCycle}};
    \node[draw=ekline,fill=white,rounded corners=2mm,below=of dut,minimum width=53mm,minimum height=8mm,align=left] (pou) {POUs · \texttt{MAIN}\\\phantom{POUs · }\texttt{FB\_Carrousel}};
    \node[draw=ekline,fill=white,rounded corners=2mm,below=of pou,minimum width=53mm,minimum height=8mm,align=left] (task) {PlcTask · exécution cyclique};
    \draw[flow] (root)--(plc);
    \draw[flow] (plc)--(dut);
    \draw[flow] (dut)--(pou);
    \draw[flow] (pou)--(task);
  \end{tikzpicture}

  \begin{primarycard}
    \textbf{Répartition des rôles}\par\small
    \texttt{MAIN} relie les variables d'entrée et de sortie. Toute la logique machine est regroupée dans \texttt{FB\_Carrousel}, ce qui facilite la lecture en ligne et la réutilisation du bloc.
  \end{primarycard}
}{
  \subsection{Fonctions présentes dans la version réalisée}
  \begin{tabularx}{\linewidth}{@{}L{37mm}Y@{}}
    \toprule
    \textbf{Fonction} & \textbf{Comportement implémenté}\\
    \midrule
    Référencement & attente du capteur inductif, puis affectation du bac 1\\
    Comptage & front montant du signal du capteur laser avec \texttt{R\_TRIG}\\
    Mémoire & huit compteurs dans \texttt{aPieceBac[1..8]}\\
    Indexation & quatre cycles complets de sortie et de rentrée du vérin\\
    Temporisation & temps de mouvement réglable par \texttt{tVerin}\\
    Bouclage & retour logique du bac 8 vers le bac 1\\
    Bac suivant déjà plein & attente de l'opérateur dans \texttt{ATTENTE\_BAC} jusqu'à la réinitialisation du bac\\
    Réinitialisation & appui bref : bac concerné ; maintien 10 s : huit compteurs\\
    Défaut & paramètres invalides, commandes logiques à \texttt{FALSE}, reprise par \texttt{INIT}\\
    \bottomrule
  \end{tabularx}

  \begin{notebox}[colback=ekwarning,colframe=orange!35]{Interprétation retenue}
    Le terme « 4 impulsions » a été traduit par quatre cycles complets du vérin : sortie temporisée, rentrée temporisée, puis incrément de \texttt{nImp}.
  \end{notebox}
}

\newpage

% ==============================================================
